\section{Protegendo seu Bolso e sua Mente \faCreditCard}

A tecnologia avança, mas o alvo do criminoso continua sendo o seu dinheiro. Aqui estão as duas defesas finais para a sua rotina:

\subsection*{1. O Cartão Bancário que Desaparece (Cartão Virtual)}
Nunca digite o número do seu cartão de plástico na internet. No aplicativo do seu banco, gere um \textbf{Cartão Virtual}. 
Ele funciona como uma nota de dinheiro mágica: você paga e o número deixa de existir. Se o site for falso e invadido, o hacker roubará um número inútil, e o seu cartão físico na sua carteira continuará seguro.

\subsection*{2. O Golpe da Maquininha no Ônibus}
Para evitar que criminosos encostem máquinas de cartão na sua bolsa no transporte público: entre no aplicativo do seu banco e \textbf{DESATIVE} a opção de \textit{\enquote{Pagamento por aproximação (NFC)}} do seu cartão físico. Se for pagar por aproximação, prefira usar o celular.

\vspace{0.5cm}

\begin{tcolorbox}[colback=NordBackground, colframe=red, title=\faBrain\ \textbf{A Regra de Ouro: Engenharia Social}]
	\textcolor{white}{\textbf{O bandido não hackeia o seu celular, ele hackeia a sua mente.}}\\
	\textcolor{white}{Se uma mensagem no WhatsApp causar \textbf{URGÊNCIA} (\enquote{sua conta será bloqueada em 30 min!}) ou \textbf{MEDO} (\enquote{mãe, perdi meu celular, manda um pix.}), \textbf{PARE E RESPIRE}. A pressa é a principal arma deles. Nunca clique em links. Feche a mensagem e ligue para a pessoa para ouvir a voz dela.}
\end{tcolorbox}